% =====================================================================
%  CAPÍTULO II — PERCURSO METODOLÓGICO
%
%  Numeração 2.1 a 2.3.2 dada pela FATEC-ID (caderno de delineamentos e
%  Manual, item 4.1.2.8). As seções seguintes (implementação, plano
%  experimental, obstáculos) detalham os procedimentos.
%
%  Varredura do autor em 2026-10-15 (lista de pedidos colada na conversa):
%  Wazlawick citado em 2.1.1 e 2.1.2 (a definição vem do PTG, escrito pelos
%  autores); critério real dos ativos (qualidade de dado, não volatilidade);
%  desdobramento explicado com as datas e os fatores; logaritmo justificado
%  com Thomaz (2019, p. 28-29); GARCH(1,1) justificado com Xing et al. (2024,
%  p. 5 do PDF); janela deslizante comparada com Xing et al. (2024) e limitada
%  pelo histórico do testador; as formas de usar o GARCH vieram do Cap. I
%  para cá, como decisão de projeto. Em 21/10/2026, a pedido do autor, a seção
%  de implementação foi dividida em subseções (sem quadro), o conceito de filtro e
%  volatility targeting voltou ao Cap. I (Harvey et al., 2018) e aqui ficou só
%  como foi implementado, mais a subseção de parâmetros;
%  validação reescrita (por que AR(2), ARMA(1,1) e
%  GARCH(1,1), tamanho das séries, semente); plano experimental com um quadro de
%  períodos e menos detalhe; sem "equação (x.y)" e com menos remissões internas.
%  Versão anterior em ia/rascunhos/antes_varredura_autor_2026-10-15/.
%
%  Em 2026-10-18, a pedido do autor, saiu a seção "Validação das
%  implementações" (testes do código em séries sintéticas, ADF em séries
%  artificiais e diagnóstico ADF/Ljung-Box na PETR4): foge do escopo
%  apresentado; pode voltar depois. Texto guardado em
%  ia/rascunhos/antes_tirar_validacao_2026-10-18/.
%
%  % REFERÊNCIA: marcações assim indicam onde uma citação deve entrar
%  depois que o autor ler a obra. Nenhuma foi inserida sem leitura.
%  Thomaz (2019), já lido pelo autor, sustenta: log-retorno (p. 28-29), AIC
%  (p. 40), máxima verossimilhança do GARCH (p. 43), ADF (p. 44) e Ljung-Box
%  (p. 47). Wazlawick (2014): definições usadas no PTG dos autores.
%
%  Procedência dos números deste capítulo:
%   - 5.881 pregões, 02/01/2003 a 21/09/2026: dados/series/PETR4_OFICIAL.csv
%   - desdobramentos (01/09/2005, 1 para 4; 28/04/2008, 1 para 2; quedas de
%     75,2% e 49,5% no COTAHIST; preços divididos por 2 e por 8):
%     dados/FONTES.md, itens 1 e 2
%   - critério dos outros ativos: dados/FONTES.md, item 9
%   - 312 configurações: dados/resultados/rodada3_b3_oficial/grades/*.txt, menos
%     as 32 do intervalo de confiança (refeito_2026-10-05/grades/r3_ic.txt), que saiu do escopo em
%     21/10/2026 a pedido do autor (INCERTEZAS item 71)
%   - janelas de otimização e teste, critérios, configurações fixas e
%     cenário de Almeida: dados/resultados/refeito_2026-10-05/PROTOCOLO.md
%   - 14.976 simulações (624 configurações x 4 ativos x 6 janelas; eram 688 e
%     16.512 com o intervalo de confiança):
%     dados/resultados/refeito_2026-10-05/analise_grades.txt
%   - janela de 120/240 pregões, reseleção a cada 21: robos/ARIMA-GARCH/ARIMA_GARCH_EA.mq5
% =====================================================================
\setdiv{CAPÍTULO II}
\chapter{Percurso metodológico}

Este capítulo descreve como a pesquisa foi conduzida. Apresenta a caracterização
do estudo, a amostra, a coleta e a análise dos dados, a implementação dos
modelos, o plano experimental e os obstáculos encontrados.

\section{Caracterização da pesquisa}

\subsection{Quanto aos objetivos}

Quanto aos objetivos, a pesquisa é explicativa, com uma parte descritiva. Segundo
\citeonline{wazlawick2021}, a pesquisa explicativa busca identificar as causas de
determinados fenômenos. Neste trabalho, ela investiga por que a inclusão de um modelo de
volatilidade altera o desempenho de uma estratégia cuja previsão de direção
permanece a mesma. A parte descritiva está no registro de como cada robô se
comportou em diferentes períodos do mercado.

\subsection{Quanto ao delineamento}

Quanto ao delineamento, trata-se de uma pesquisa quase experimental. A pesquisa
experimental envolve a manipulação controlada de variáveis e a observação dos
efeitos resultantes. Quando os objetos estudados não podem ser selecionados de
forma aleatória, trata-se de um quase experimento \cite{wazlawick2021}. É o caso
deste trabalho, em que a série histórica de um ativo é única e não pode ser
sorteada.
% Wazlawick (2021), cap. 4, pesquisa experimental: "ao invés de selecionar os
% objetos aleatoriamente, eles são selecionados por conveniência [...] trata-se de
% um quase experimento" (p. 45 do PDF, sem paginação impressa). Em 19/10/2026 o
% autor pediu para não citar documentos da própria FATEC-ID.

O problema e as hipóteses foram formulados na introdução. As variáveis são os parâmetros
configuráveis do robô, e o plano experimental foi definido antes da observação
dos resultados. Os sujeitos são as séries históricas dos ativos, e o ambiente é a
plataforma de simulação.

\section{Caracterização do lugar e da amostra de pesquisa}

O lugar da pesquisa é computacional. Os experimentos foram conduzidos no testador
de estratégias do MetaTrader 5, que simula a execução de ordens sobre séries
históricas com as mesmas regras de negociação que usaria em tempo real. A
plataforma foi escolhida, em primeiro lugar, por ser a mesma utilizada no estudo
que este trabalho continua \cite{almeida2016}, o que permitiu reproduzir a
estratégia de referência no mesmo ambiente. Pesou também o fato de ela reunir em
um só lugar a programação dos robôs, o histórico de preços e o teste sobre dados
passados, e de contar com uma grande comunidade de usuários e desenvolvedores,
motivos que levaram \citeonline{batista2024} à mesma escolha. Ela é nativa do sistema Windows, e os experimentos
apresentados aqui foram executados em ambiente Linux, por meio da camada de
compatibilidade Wine, pelos motivos expostos ao final deste capítulo.

A amostra é composta por séries históricas diárias de ações negociadas na B3. O
ativo principal é a PETR4, ação preferencial da Petrobras, escolhida por três
motivos. É o ativo estudado por \citeonline{almeida2016}, o que permite a
comparação direta. Apresenta volatilidade acentuada, característica que dá sentido
à aplicação de um modelo de variância condicional. E possui série longa e contínua
no arquivo histórico da bolsa, com 5.881 pregões, de 2 de janeiro de 2003 a 21 de
setembro de 2026.

Para verificar se os resultados se repetem fora do ativo principal, foram
incluídas outras três ações, VALE3, da mineração, USIM5, da siderurgia, e POMO4,
de bens de capital. O critério de escolha, definido antes da observação de
qualquer resultado, considerou a qualidade dos dados e a diversidade de setores.
Cada ação precisava ter série completa de 2003 a 2026 no arquivo da bolsa, sem
mudanças na forma de cotação que impedissem a reconstrução dos preços, e preço
sempre acima de R\$ 3, para que o arredondamento em centavos não distorcesse os
retornos. As ações também deveriam pertencer a setores diferentes entre si.

O recorte temporal foi dividido em períodos, e não tratado como um período único,
porque o comportamento de uma estratégia depende do momento do mercado. Foram
examinados um período de crise, de 2007 a 2011, um de queda prolongada seguida de
recuperação, de 2013 a 2016, o intervalo de 2017 a 2024, que inclui a crise da
pandemia, o período mais recente, de 2025 a 2026, e o período mais longo
disponível, de 2005 a 2026.

\section{Procedimentos para coleta e análise de dados}

\subsection{Técnicas para coleta de dados}

Os dados foram obtidos de fontes públicas oficiais, e não do servidor da
corretora, que disponibiliza cotações apenas a partir de setembro de 2021 e pode
alterar o histórico sem aviso. As cotações vieram do arquivo COTAHIST, série
histórica publicada pela própria B3, que registra, para cada pregão e cada ação,
os preços de abertura, máximo, mínimo e fechamento e o volume negociado
\cite{b3Cotahist}. A taxa Selic diária foi obtida da série 11 do Sistema Gerenciador de Séries
Temporais do Banco Central \cite{bcbSgs11}.

Sobre as cotações foi aplicada uma única transformação, a correção por
desdobramentos. No desdobramento, cada ação é dividida em várias, e o preço cai na
mesma proporção, sem que o investidor perca nada, já que passa a ter mais ações.
Na PETR4 houve dois desdobramentos no período. Em 1º de setembro de 2005, cada ação
virou quatro, e em 28 de abril de 2008, cada ação virou duas. No arquivo da bolsa,
esses dias aparecem como quedas de 75,2\% e de 49,5\%, que não aconteceram de fato
e que o modelo interpretaria como dois tombos do mercado. Para corrigir, os preços
anteriores a 28 de abril de 2008 foram divididos por dois, e os anteriores a 1º de
setembro de 2005, por oito, isto é, por quatro e depois por dois. Com isso, a série
passa a registrar apenas as variações reais de preço. Os dois eventos foram
confirmados nos eventos corporativos que a B3 divulga para a Petrobras
\cite{b3Petrobras} e, o de 2005, que esse registro não alcança, no comunicado em
que a própria companhia anunciou o desdobramento \cite{petrobras2005}.

Os proventos, isto é, os dividendos e os juros sobre capital próprio pagos aos
acionistas, não foram incorporados à série, para que os modelos trabalhassem
apenas com as variações de preço negociadas em bolsa. Eles foram obtidos à parte,
para a análise complementar do Capítulo III, no histórico de proventos em dinheiro
que a B3 divulga para a Petrobras \cite{b3Petrobras}. Para cada provento, esse
histórico informa a última data com direito a recebê-lo e o seu valor como
percentual do fechamento daquele dia. O valor por ação foi obtido aplicando esse
percentual ao fechamento da série corrigida, o que o deixa na mesma escala dos
preços corrigidos pelos desdobramentos.
% Procedência: ferramentas/scripts/proventos_b3.sh (corporateActionPrice x
% fechamento da PETR4_OFICIAL na data-com; data-ex = pregão seguinte), arquivos
% dados/fontes/b3/proventos_PETR_historico_api_b3_2026-10-21/, dados/FONTES.md item 3.

As séries preparadas foram importadas para a plataforma como símbolos
personalizados, recurso que permite ao testador operar sobre dados próprios.

\subsection{Natureza da análise de dados}

A análise é quantitativa. Os robôs foram comparados pelas métricas apresentadas
no Capítulo I. Como o GARCH decide se o robô opera e com quanto, e não a direção
da operação, que continua a cargo do ARIMA, as métricas principais são as de
risco, a queda máxima e a volatilidade. O retorno acumulado acompanha essas
métricas como resultado complementar. A taxa de acerto verifica se o GARCH
melhorou o acerto direcional, como previa a primeira hipótese, e a quantidade de
sinais vetados ajuda a explicar o funcionamento do filtro.

Para verificar se as diferenças entre os robôs eram consistentes, e não obra do
acaso, observou-se se elas se repetiam entre os períodos, entre as configurações e
entre os ativos. Se o GARCH não fizesse diferença, a versão com ele seria melhor
em cerca de metade das comparações, como no lançamento de uma moeda.

\section{Implementação dos modelos}

O MetaTrader 5 oferece, na biblioteca padrão do MQL5, indicadores de análise
técnica como médias móveis e osciladores. Os modelos ARIMA e GARCH foram
implementados para este
trabalho. A implementação se apoiou na documentação oficial da linguagem
\cite{metaquotesMql5} e nos conteúdos compartilhados pela comunidade de
desenvolvedores do MetaTrader 5, cujos exemplos e trechos de código serviram de
ponto de partida e foram adaptados às necessidades do robô. Para os cálculos
numéricos mais pesados, a resolução de sistemas lineares e a otimização com
restrições, foram usadas rotinas da ALGLIB, biblioteca numérica que acompanha a
instalação da plataforma.

O robô de médias móveis de \citeonline{almeida2016} seguiu o mesmo caminho. Ele
foi reescrito a partir do algoritmo publicado pelo autor, com o indicador de média
móvel da própria plataforma, e compartilha com o ARIMA-GARCH a rotina de envio de
ordens e de medição dos resultados, para que os dois sejam avaliados nas mesmas
condições.

As subseções a seguir descrevem a implementação do ARIMA-GARCH parte por parte,
da preparação dos preços aos parâmetros do robô.
% Frase genérica, sem referência, a pedido do autor (21/10/2026). Exemplos do que
% ele leu, se a banca perguntar: Dube (2023), ARIMA em MQL5 pelo método de Powell,
% https://www.mql5.com/pt/articles/12583; Dube (2024), teste ADF em MQL5,
% https://www.mql5.com/pt/articles/13991. O código não foi tirado deles.

\subsection{Preparação dos preços}

Os modelos foram aplicados ao logaritmo do preço, e não ao preço em reais. A
diferença entre os logaritmos de dois preços seguidos é o retorno logarítmico, que
mede a variação percentual do dia \cite{thomaz2019}. Trabalhar com retornos tem
duas vantagens. A primeira é que eles são livres de escala, ou seja, uma alta de
1\% significa o mesmo para uma ação cotada a R\$ 5 ou a R\$ 40, de modo que os
parâmetros do robô valem para qualquer ativo e período. A segunda é que os
retornos costumam ser estacionários, ao contrário dos preços, condição exigida
pelas partes autorregressiva e de médias móveis do ARIMA \cite{thomaz2019}.

Como consequência, a volatilidade estimada pelo GARCH também sai em porcentagem.
Uma volatilidade de 0,02 indica uma oscilação típica de 2\% ao dia, enquanto, com
o preço em reais, o mesmo número significaria dois centavos e mudaria de sentido
conforme o preço da ação.

\subsection{Estimação do ARIMA}

O ARIMA foi estimado pelo método de \citeonline{hannan1982}, que divide a
estimação em duas regressões de mínimos quadrados ordinários, o método que escolhe
os coeficientes que tornam mínima a soma dos quadrados dos erros de ajuste. A dificuldade de um modelo com
termos de médias móveis é que os erros passados não podem ser observados, pois
dependem dos próprios coeficientes que se quer estimar. A primeira regressão
ajusta um modelo autorregressivo longo, cujos resíduos servem de estimativa desses
erros, e a segunda usa essas estimativas para obter os coeficientes.

A escolha do método foi uma decisão de projeto. Ele não exige taxa de aprendizado
nem número de iterações e não depende da escala dos preços, o que o torna estável
para um modelo que é reestimado todos os dias. O método também é usado na
identificação automática de modelos ARIMA, em que muitas combinações de ordens
precisam ser estimadas, por exigir apenas regressões lineares e fornecer
estimativas próximas às de máxima verossimilhança \cite{gomez1998}.
% Hannan e Rissanen (1982) e Gómez e Maravall (1998, p. 10 e 27-28) citados em
% 20/10/2026, autorizados pelo autor. PDF em bibliografia/pdfs/Gomez-1998.pdf.
% ia/PENDENCIAS.md, item 1: pode ser trocado por máxima verossimilhança condicional
% (Thomaz, 2019, p. 42-43), partindo das estimativas de Hannan-Rissanen.

A inclusão da constante $c$ do modelo foi tratada como parâmetro da configuração.
Com $d = 1$, ela representa a tendência média dos retornos da janela. No caso mais
simples, sem termos autorregressivos nem de médias móveis, a previsão é a própria
média dos retornos da janela, e o robô passa a seguir a direção da tendência
recente. Sem a constante, a previsão depende apenas dos retornos e dos erros
passados ponderados pelos coeficientes, e o modelo sem nenhum termo foi descartado
da busca, porque a sua previsão seria sempre zero. Com $d = 0$, a constante foi
sempre incluída, porque o logaritmo do preço não oscila em torno de zero.
% Código: CArima::Definir (constante obrigatória com d = 0), HannanRissanen
% (p = q = 0: constante = média da série diferenciada) e BuscarMelhorPQ
% (ARIMA(0,d,0) sem constante é pulado).

Os coeficientes saem das regressões, mas a comparação entre as ordens, feita pelo
AIC, exige a verossimilhança de cada modelo. Ela foi calculada na forma
condicional, supondo que os erros sigam a distribuição normal, caso em que depende
da soma dos quadrados dos resíduos do modelo \cite[p.~42]{thomaz2019}. Quanto
menores os erros que o modelo comete sobre a janela, maior a sua verossimilhança.
Todas as combinações candidatas foram avaliadas sobre os mesmos pregões da janela,
para que nenhuma levasse vantagem por somar mais ou menos dias.
% Código: CArima::CalcularVerossimilhanca (logL = -n/2 [ln(2 pi) + ln(sigma2) + 1],
% sigma2 = SQR/n) e DefinirAquecimentoMinimo chamado com o mesmo valor em todos os
% candidatos em BuscarMelhorPQ.

As ordens do modelo foram escolhidas pelo próprio robô. A ordem de integração $d$
foi determinada pela aplicação sequencial do teste ADF \cite{thomaz2019}, ao nível
de significância de 5\%. Testa-se a série e, enquanto a hipótese de raiz unitária
não for rejeitada, testa-se a série diferenciada mais uma vez. Com $d = 1$, o ARIMA modela os retornos. Com $d = 0$,
escolhido quando o próprio logaritmo do preço se mostra estacionário dentro da
janela, ele modela o logaritmo do preço diretamente. As ordens $p$ e $q$ foram
escolhidas por busca em grade, de zero a quatro, pelo menor AIC
\cite{thomaz2019}, procedimento semelhante ao de \citeonline{xing2024}. As ordens
são reavaliadas a cada 21 pregões, cerca de um mês, de modo que podem mudar ao
longo de um teste.

\subsection{Estimação do GARCH}

O GARCH foi estimado sobre os resíduos do ARIMA, isto é, sobre a parte do retorno
que o ARIMA não conseguiu prever. Foi adotada a especificação GARCH(1,1), que exige
poucos parâmetros e é a escolha predominante nos estudos de previsão de ações
\cite{xing2024}. Estimar o modelo significa encontrar os valores de $\omega$,
$\alpha$ e $\beta$ que melhor descrevem como a agitação desses resíduos variou ao
longo da janela.

Para isso foi usada a máxima verossimilhança, que escolhe os valores dos
parâmetros que tornam os dados observados os mais prováveis segundo o modelo
\cite[p.~41]{thomaz2019}. Na prática, cada combinação candidata é posta à prova
sobre a janela. Com ela, a equação do GARCH apresentada no Capítulo I indica, para
cada dia, a variância que o modelo esperava, e essa expectativa é comparada com o
erro que de fato ocorreu. Supondo que os erros sigam a distribuição normal, um erro
grande em um dia para o qual o modelo esperava calma é pouco provável e reduz a
verossimilhança, enquanto o mesmo erro em um dia para o qual o modelo já esperava
agitação é plausível \cite[p.~43]{thomaz2019}. A melhor combinação é a que
antecipa melhor, ao longo de toda a janela, quais dias seriam agitados e quais
seriam calmos.

Ao contrário do que ocorre no ARIMA, essa busca não se resolve com uma regressão.
Foi usado um otimizador numérico, que experimenta valores dos parâmetros, verifica
se a verossimilhança aumentou e repete o processo até não conseguir mais
melhorá-la. As restrições apresentadas no Capítulo I foram impostas dentro do
próprio otimizador, de modo que ele só experimenta valores válidos.

Um cuidado foi necessário com $\omega$. Ele é a parte fixa da variância de cada
dia, que não depende dos choques nem da variância anterior, e por isso fica na
escala da variância diária dos retornos, um número muito pequeno, enquanto $\alpha$
e $\beta$ ficam entre zero e um. Para que o otimizador não precisasse lidar com
números de tamanhos tão diferentes, $\omega$ foi expresso como uma fração da
variância dos resíduos da janela e convertido de volta ao final. Concluída a
estimação, o valor usado na decisão é a volatilidade prevista para o pregão
seguinte.
% Código: CGarch::Ajustar (MinBLEIC da ALGLIB, gradiente por diferenças finitas;
% omega relativo = omega / variância amostral dos resíduos, convertido de volta em
% m_omega = x_final[0] * m_variancia_amostral) e a função objetivo, que soma a
% log-densidade normal de cada resíduo com a variância daquele dia.

\subsection{Janela de estimação}

Os dois modelos são reestimados a cada pregão sobre uma janela deslizante, isto é,
sobre os últimos 120 ou 240 pregões já encerrados, conforme a configuração. A cada
dia, o pregão mais antigo sai da janela e o mais recente entra, o que permite ao
modelo acompanhar as mudanças de comportamento do mercado. A janela nunca inclui o
pregão em andamento, para que a decisão do dia não use uma informação que só
estaria disponível no seu encerramento.

\citeonline{xing2024} compararam esse tipo de janela com a janela expansiva, que
acumula todo o histórico desde o início, e no seu estudo o ARIMA teve melhor
desempenho com a janela expansiva. A janela expansiva não foi testada neste
trabalho. O tamanho máximo da janela, de 240 pregões, decorreu da forma como os
testes foram montados. Cada teste foi iniciado no testador na data inicial do
período, e observou-se que, nessa condição, o histórico disponível antes do início
começa em 1º de janeiro do ano anterior e soma cerca de 250 pregões. Janelas
maiores exigiriam iniciar a simulação antes do período e só começar a operar na
data inicial, alternativa que não foi adotada neste trabalho.
% 2026-10-29: a pedido do autor, a frase deixou de dizer que o testador impede
% janelas maiores. A documentação da MetaQuotes só diz que o testador baixa dados
% para formar no mínimo 100 barras antes do início; o 1º de janeiro é observação
% (INCERTEZAS item 8). O estudo de 26/10 (janelas_240_500_2026-10-26) usou 500
% começando o testador dois anos antes, com Inicio_Operacao, e não entra no TG.

\subsection{Funcionamento do robô}

Os dois modelos foram reunidos em um único \textit{Expert Advisor}, com um
parâmetro que liga e desliga o GARCH. Com ele desligado, o programa percorre o
mesmo caminho de previsão, a mesma janela de estimação e a mesma rotina de envio
de ordens, de modo que qualquer diferença entre as duas execuções só pode vir do
GARCH.

A divisão de papéis é estrita. O ARIMA é o único responsável por indicar o lado da
operação, e o GARCH nunca gera sinal de direção. A cada novo pregão, o robô
carrega a janela, ajusta o ARIMA, obtém a previsão do retorno do dia, ajusta o
GARCH sobre os resíduos e obtém a volatilidade prevista. Se o GARCH não vetar a
operação, o sinal da previsão do ARIMA determina a direção, e a ordem é executada
na abertura do pregão.

\subsection{Uso do GARCH na decisão}

As duas maneiras de usar a volatilidade prevista apresentadas no Capítulo I, o
filtro e o ajuste da posição, foram implementadas no robô, cada uma sozinha e as
duas combinadas.

No filtro, o limiar pode ser definido de duas formas. Na primeira, ele é um valor
fixo de volatilidade diária, como 2\% ao dia. Na segunda, ele é um percentil da
volatilidade que o próprio GARCH estimou dentro da janela. Com o percentil 70, por
exemplo, a operação é vetada quando a volatilidade prevista para o dia supera 70\%
dos valores estimados na janela. O limiar por percentil se adapta ao ativo e ao
período, porque o que é um dia agitado muda com o tempo.

No ajuste da posição, o tamanho da posição é multiplicado pela razão entre o alvo
e a volatilidade prevista, dentro de limites mínimo e máximo. Com um alvo de 1\% ao
dia e uma volatilidade prevista de 2\%, o robô compra metade do tamanho normal.

O ajuste da posição sozinho foi testado, mas não foi escolhido em nenhum período
nem em nenhuma ação. Todos os robôs escolhidos usam o filtro, sozinho ou combinado
com o ajuste da posição.

Nas formas com filtro, resta decidir o que fazer com a posição que já está aberta.
Por padrão, o veto só impede novas compras, e o robô continua com o que já tinha
comprado. Uma opção do robô permite que o veto também encerre a posição, retirando
o capital do mercado enquanto a volatilidade prevista estiver alta. O ajuste do
tamanho, por sua vez, vale apenas para as novas compras.

\subsection{Parâmetros do robô}
\label{sec:parametros}

O comportamento do robô é definido por parâmetros de entrada, que o MetaTrader 5
permite alterar sem modificar o código. Eles se dividem em três grupos.

Os parâmetros do ARIMA são o tamanho da janela de estimação, a inclusão ou não de
uma constante no modelo e uma previsão mínima, abaixo da qual o robô não opera
mesmo que a previsão indique alta. Os parâmetros do GARCH são a chave que liga e
desliga o modelo, a forma de uso, o tipo e o valor do limiar, o alvo de
volatilidade e a opção de encerrar a posição no veto. Os parâmetros de execução são a fração do patrimônio aplicada em cada
operação e a restrição de operar apenas na ponta compradora.

Alguns parâmetros ficaram fixos em todos os testes, como a ordem máxima de zero a
quatro do ARIMA e a reavaliação das ordens a cada 21 pregões. Os demais foram
variados na otimização, com os valores apresentados no plano experimental.

\section{Plano experimental}

O plano experimental foi definido antes da observação dos resultados. Esta seção
descreve, em partes, as condições das simulações, a forma de medir o resultado, a
reprodução do estudo de referência, as configurações testadas, a forma de escolher
entre elas, os robôs comparados e os testes complementares.

\subsection{Condições das simulações}

Todas as simulações usaram os preços de abertura de cada pregão, já que o robô
decide uma vez por dia e executa na abertura. No cenário principal, partiram de um
capital de R\$ 100 mil, aplicaram a cada operação 95\% do patrimônio e operaram
apenas na ponta compradora. A venda a descoberto, em que o robô vende uma ação que
não tem para recomprá-la mais barato, ficou fora do cenário principal. Nela, o
robô passa a apostar também na queda, e o resultado depende ainda mais da direção
prevista pelo ARIMA, enquanto o foco deste trabalho é a proteção dada pelo GARCH.
O efeito dela foi medido à parte, nos testes complementares. A reprodução do
estudo de referência usou as condições do autor, descritas mais adiante.

\subsection{Medição dos resultados}

O resultado de cada simulação foi medido pelo patrimônio da conta simulada pelo
testador, sem nenhum ajuste, chamado aqui de resultado bruto. A comparação entre
os modelos é feita sobre ele, porque depende apenas das decisões de cada robô. A
compra e manutenção da ação foi medida pela variação do preço.

Na escolha das configurações, descrita adiante, o desempenho foi medido de forma
mais próxima da operação real, com dois ajustes. O primeiro desconta, a cada ordem,
0,03\% do valor negociado, que é a soma das tarifas de negociação, de contraparte
central e de transferência que a B3 cobra de cada investidor, na compra e na venda,
em operações com ações à vista que não são \textit{day trade}, na faixa de quem
negocia em média até R\$ 3 milhões por dia \cite{b3Tarifas}. A corretagem foi
considerada zero, e a tarifa, que é a vigente em 2026, foi aplicada a todos os
períodos. O segundo acrescenta a remuneração pela taxa Selic diária do capital que
fica fora do mercado, para que configurações que passam tempos diferentes fora
dele sejam comparadas de forma justa.
% Custo de 0,03% por lado (input Custo_Por_Lado do EA): tabela de tarifas da B3,
% à vista, sem day trade, ADTV mensal até R$ 3 milhões: negociação 0,00500% + CCP
% 0,02240% + TTA 0,0026% = 0,0300% (dados/fontes/b3/tarifas_acoes_a_vista_b3_2026-10-21.html,
% dados/FONTES.md item 8).

Como análise complementar, o resultado foi medido também com os dois ajustes e
com os proventos. Cada provento foi creditado no primeiro pregão em que a ação
passa a ser negociada sem direito a ele, na proporção do patrimônio que o robô
mantinha aplicado na ação naquele dia. Na compra e manutenção, o provento foi
somado ao preço do mesmo dia. O resultado é uma estimativa do que o investidor
teria recebido, porque não considera impostos nem o prazo entre essa data e o
pagamento.
% Desde 27/10/2026 (pedido do autor, depois da leitura da orientadora) os
% proventos voltam, só com a base da B3 (camada 4 do Desempenho.mqh:
% exposicao x provento / preço anterior, na data-ex; compra e manutenção:
% (P + provento) / P anterior). Resultados em
% dados/resultados/refeito_2026-10-05/proventos_b3/. A base do Yahoo não entra.

Além das medidas de desempenho, o robô pode registrar em arquivo, a cada pregão,
as ordens do ARIMA em uso, os parâmetros estimados, a volatilidade prevista, o
limiar e a decisão tomada. Esse registro não interfere nas decisões e foi usado
para acompanhar o funcionamento dos modelos, como no episódio de 2020 descrito
no Capítulo III.

\subsection{Reprodução do estudo de referência}

Antes das comparações, a estratégia de médias móveis reproduzida foi submetida aos
testes de \citeonline{almeida2016}, para garantir que a comparação fosse feita com
a estratégia dele. Os testes foram repetidos com as mesmas datas, o mesmo capital
de R\$ 50 mil, as mesmas médias de 10 e 20 períodos, sem custos e com operações de
compra e venda. O autor declara mil ações para o teste de 2013 a 2016, mas não
informa o volume dos testes de crise, e, como ele menciona cem ações como volume
mínimo, esse foi o volume adotado nesses testes.

A reprodução foi feita de duas formas. A primeira usou a série de preços deste
trabalho, a mesma dos demais experimentos. A segunda testou uma hipótese sobre o
ambiente do autor, que usou o histórico da corretora sem descrever como ele era.
A hipótese é que o histórico da PETR4 no servidor da corretora começasse em 2008.
O testador precisa de pelo menos cem pregões de histórico antes da data inicial
e, quando não os encontra, adia o início do teste por conta própria
\cite{metaquotesTeste}. Assim, os testes de crise do autor teriam começado, na
prática, em junho de 2008. Para testar essa hipótese, a série deste trabalho foi
cortada para começar em 2 de janeiro de 2008, sem nenhuma outra alteração nos
preços. Trata-se de uma hipótese, e não de um fato, porque o servidor usado pelo
autor não existe mais e não há como conferir o histórico que ele baixou.
% 2026-10-29, a pedido do autor: a segunda forma passou a ser só o histórico desde
% 2008, sem proventos e sem arredondar as médias (a combinação que chega mais perto
% dos números do autor). Série PETR4_DESDE2008 (series_replica_almeida.sh); adiamento
% no log inicio_adiado_testador.log ("start time changed to 2008.05.31").
% Versão anterior, com proventos e arredondamento:
% ia/rascunhos/antes_objetivos_risco_2026-10-29/cap2-metodologia_antes_replica.tex.

\subsection{Configurações testadas}

As configurações testadas foram geradas a partir dos parâmetros descritos na
Seção~\ref{sec:parametros}, pelo otimizador do testador de estratégias do
MetaTrader 5, no modo de busca completa, em que a própria plataforma executa
uma simulação para cada combinação dos valores informados.

Três parâmetros variaram em todas as formas de uso do GARCH. O primeiro foi a
inclusão ou não de uma constante no ARIMA. O segundo, a janela de estimação, de 120
ou de 240 pregões. O terceiro, uma previsão mínima, de zero ou de 0,1\%, abaixo da
qual o robô não opera. Essas três escolhas formam oito combinações, e cada forma de
uso do GARCH acrescentou a elas os seus próprios parâmetros e, quando havia veto, a
opção de encerrar ou não a posição, como mostra o Quadro~\ref{qua:grade}. Ao todo,
foram 304 configurações com o GARCH e oito sem ele, 312 no total.

% Procedência: dados/resultados/refeito_2026-10-05/grades/r3_*.txt (valores
% início||passo||fim de cada input; Optimization=1 em ferramentas/scripts/otimizar.sh)
\begin{quadro}[htbp]
\centering
\caption{Configurações testadas na otimização}
\label{qua:grade}
\fontsize{10}{12}\selectfont
\begin{tabular}{|>{\raggedright\arraybackslash}p{4.2cm}|>{\raggedright\arraybackslash}p{7.0cm}|>{\centering\arraybackslash}p{3.2cm}|}
\hline
Forma de uso do GARCH & Parâmetros próprios e valores testados & Configurações \\
\hline
Sem GARCH (ARIMA) & nenhum & 8 \\
\hline
Filtro, limiar fixo & limiar de 2\%, 3\% ou 4\% ao dia, encerrar ou não a posição no veto & 8 $\times$ 3 $\times$ 2 = 48 \\
\hline
Filtro, limiar por percentil & percentil 50, 70 ou 90, encerrar ou não a posição no veto & 8 $\times$ 3 $\times$ 2 = 48 \\
\hline
Filtro e ajuste da posição, limiar fixo & limiar de 2\%, 3\% ou 4\% ao dia, alvo de 1\% ou 2\% ao dia, encerrar ou não a posição no veto & 8 $\times$ 3 $\times$ 2 $\times$ 2 = 96 \\
\hline
Filtro e ajuste da posição, limiar por percentil & percentil 50, 70 ou 90, alvo de 1\% ou 2\% ao dia, encerrar ou não a posição no veto & 8 $\times$ 3 $\times$ 2 $\times$ 2 = 96 \\
\hline
Ajuste da posição & alvo de 1\% ou 2\% ao dia & 8 $\times$ 2 = 16 \\
\hline
Total & & 312 \\
\hline
\end{tabular}
\fonte{Elaboração própria.}
\nota{O número 8 corresponde às combinações comuns a todas as formas: constante no ARIMA (sim ou não), janela de estimação (120 ou 240 pregões) e previsão mínima para operar (zero ou 0,1\%).}
\end{quadro}

Cada forma de uso foi otimizada separadamente, para que parâmetros que não se
aplicam a ela não gerassem simulações repetidas. Em cada simulação, o próprio
robô gravou em arquivo as medidas de desempenho, e o critério de escolha foi
aplicado sobre esses arquivos.

\subsection{Escolha da configuração}

Escolher a melhor das 312 configurações no mesmo período em que o desempenho é
medido daria uma vantagem falsa, porque, olhando para trás, sempre existe alguma
combinação que teria dado certo. Por isso, a escolha seguiu o princípio do teste
avançado, ou \textit{forward}, do testador de estratégias, em três passos. No primeiro, as 312
configurações foram simuladas nos três anos anteriores ao período de teste. No
segundo, a lista de resultados foi ordenada por um critério definido antes, e a
primeira configuração da lista foi a escolhida. No terceiro, só a configuração
escolhida foi executada no período de teste, que não tinha sido usado na escolha,
sem nenhuma alteração. Durante o teste, os coeficientes do ARIMA e do GARCH
continuam a ser reestimados a cada pregão, como descrito na seção anterior. O que
fica fixo é a configuração.

Para cada período de teste, os três passos foram refeitos do zero, com os três
anos anteriores a ele. Por isso, a configuração escolhida, inclusive a forma de
uso do GARCH, pode mudar de um período para outro. O Quadro~\ref{qua:periodos}
mostra os períodos usados.

\begin{quadro}[htbp]
\centering
\caption{Períodos de otimização e de teste}
\label{qua:periodos}
\fontsize{10}{12}\selectfont
\begin{tabular}{|>{\raggedright\arraybackslash}p{3.4cm}|>{\raggedright\arraybackslash}p{3.8cm}|>{\raggedright\arraybackslash}p{7.2cm}|}
\hline
Otimização & Teste & Característica do período de teste \\
\hline
2004 a 2006 & 2007 a 2011 & crise financeira de 2008 \\
\hline
2010 a 2012 & 2013 a 2016 & queda prolongada da PETR4 e recuperação em 2016 \\
\hline
2014 a 2016 & 2017 a 2024 & teste principal, por ser o mais longo, inclui a crise da pandemia em 2020 \\
\hline
2022 a 2024 & jan. 2025 a set. 2026 & período mais recente disponível \\
\hline
2010 a 2012 & 2005 a 2026 e 2017 a 2024 & configurações fixas, mantidas sem nova escolha \\
\hline
\end{tabular}
\fonte{Elaboração própria.}
\end{quadro}

Foram usados dois critérios para ordenar a lista, porque não existe uma única
forma de definir a melhor configuração. O ARIMA-GARCH Sharpe é a primeira
configuração da lista ordenada pelo índice de Sharpe, ou seja, a de maior ganho
por unidade de oscilação na otimização. O ARIMA-GARCH Calmar é a primeira da lista
ordenada pelo índice de Calmar, a de maior ganho por unidade de queda máxima. O
nome indica apenas o critério usado para ordenar a lista, e não o resultado obtido
no teste. Na otimização, os dois índices foram calculados com os dois ajustes
descritos na medição dos resultados, e o de Sharpe, sobre o ganho acima da Selic.

\subsection{Robôs comparados}

Em cada período, as versões com GARCH foram comparadas a três referências. A
primeira é o melhor robô sem o GARCH, escolhido entre as oito configurações sem o
modelo pelo mesmo critério do ARIMA-GARCH Sharpe, e que aparece nas tabelas
simplesmente como ARIMA. A segunda é a estratégia de médias móveis de
\citeonline{almeida2016}, reproduzida neste trabalho e executada nos mesmos
períodos e nas mesmas condições, inclusive apenas na ponta compradora, ao
contrário do teste original, que também vendia a descoberto. A terceira é a
compra e manutenção da própria PETR4.

Essas comparações colocam lado a lado robôs com configurações diferentes. Para
confirmar que as diferenças vinham do GARCH, e não dos demais parâmetros, cada
versão com o GARCH também foi executada com o modelo desligado e todo o resto
igual. Como o robô tem um único parâmetro que liga e desliga o GARCH, essa
comparação isola o efeito dele.

\subsection{Testes complementares}

Dois testes complementares ampliaram a verificação. No primeiro, o ARIMA-GARCH
Sharpe e o ARIMA-GARCH Calmar escolhidos em 2010 a 2012 foram mantidos fixos, sem
refazer a escolha, assim como \citeonline{almeida2016} manteve fixas as suas
médias de 10 e 20 períodos. Eles foram testados em 2017 a 2024, período
inteiramente posterior à escolha, e em 2005 a 2026, que inclui os anos da escolha
e por isso não é um teste fora da amostra, mas mostra se o comportamento se mantém
no prazo mais longo disponível.

No segundo, o mesmo procedimento foi repetido na PETR4, na VALE3, na USIM5 e na
POMO4, nos períodos de 2007 a 2011, 2013 a 2016 e 2025 a 2026, e cada configuração
com GARCH foi comparada à mesma configuração sem ele. Nesse teste, a grade incluiu
também a opção de vender a descoberto, o que elevou o número de configurações para
624 e o total para 14.976 simulações. Como no cenário principal, a escolha
considerou apenas as configurações que operam na ponta compradora, e as demais
serviram para medir o efeito da venda a descoberto.

\section{Obstáculos encontrados}

O desenvolvimento enfrentou dois obstáculos que alteraram o cronograma original e
influenciaram decisões metodológicas.

O primeiro foi a perda do ambiente Windows. O desenvolvimento começou em uma
máquina com Windows e MetaTrader 5 instalado nativamente, e a perda desse
equipamento obrigou a migração para Linux, onde a plataforma não tem suporte
oficial e precisa ser executada por camada de compatibilidade. A solução funcionou,
mas trouxe instabilidades ocasionais e exigiu que a automação dos experimentos
fosse refeita em linha de comando.

O segundo foi o acesso aos dados históricos pela corretora. A corretora utilizada
inicialmente, a Rico, deixou de oferecer suporte ao MetaTrader 5, o que exigiu a
abertura de conta em outra instituição. A segunda corretora, a Clear, oferece a
plataforma, mas disponibiliza cotações apenas a partir de setembro de 2021, o que
não atendia aos períodos de teste, que retrocedem a 2003. Esse obstáculo levou à
adoção do arquivo oficial da bolsa como fonte, o que tornou a série pública,
verificável e livre de alterações posteriores.
